\section{Introduction}
\label{sec:v7-introduction}
What can be learned from a rare collision record, and how many independent preparations can be used before a limiting description becomes distinguishable from the physical experiment? These are different questions. An exact probability coefficient can identify a parameter without providing a feasible acquisition procedure. A limiting conditional distribution can describe each fixed observation accurately without approximating an experiment whose sample size increases. We study both questions in a setting where the probability law, the selection event, and the preparation cost can be derived from the underlying dynamics.

The observation is a collision record generated by a unit-speed billiard with an initial condition drawn from normalized phase volume. Close to the earliest time permitting a given number of collisions, the relevant trajectory alternates between two facing contacts. Increasing the number of flights makes the selected event exponentially rare. It does not increase the dimension of the free endpoint variables. The resulting combination of a two-dimensional transverse action and a shrinking physical flux is the source of the statistical problem.

Our central analytical result is a uniform relative two-boundary law. It separates the nonlinear stationary action into two half-line actions and the normalized physical flux into their determinant amplitudes, on a neighborhood independent of the number of flights. The word relative is essential: the reference mixed derivative is exponentially small, even though the endpoint Hessian remains uniformly positive. An absolute approximation to the action or flux would not control the normalized probability experiment. The proof establishes summable endpoint localization and trace-norm control of a relative determinant before integrating the physical residual-time coordinate.

This construction has an operational implication. Let $j$ be the number of complete flights, let $d>0$ be the excess over its channel onset, and let $\gamma>0$ be the channel instability parameter. The probability $p_{j,d}$ of a selected record is of order $d^2e^{-j\gamma}$. For the observed endpoints and the first residual time, the conditional finite and boundary laws differ in total variation by at most $C\tau^j$, with $0<\tau<1$. Restoring the failure outcome gives an error at most $Cp_{j,d}\tau^j$ per preparation. Thus a sample of $n$ preparations can be replaced by the boundary model whenever its expected number of successes, $np_{j,d}$, times $\tau^j$ tends to zero. This is a simultaneous statement about record length and experiment size, not a succession of fixed-$j$ expansions.

We determine a sharp benchmark for this replacement in the identical-even contact class. Put $u_j=e^{-j\gamma}$ and
\[
 p_j^0(d)=\frac{d^2}{2A\sinh(j\gamma)}.
\]
The quadratic successful experiments have exact one-record distance $(2/\pi)\arcsin u_j$. Their product distance is one minus the corresponding overlap probability raised to the sample size. More importantly, uniform nonlinear estimates pass this calculation to actual positive-offset physical laws. If $d_j=o(u_j)$ and $k_ju_j\to b$, the distance between the two experiments with $k_j$ successful records tends to
\begin{equation}\label{eq:v7-intro-phase}
                  1-\exp(-2b/\pi).
\end{equation}
The same limit holds for raw preparations when $n_jp_j^0(d_j)u_j\to b$. We include $b=0$ and $b=\infty$. Consequently the boundary model is indistinguishable at a smaller scale, has a nontrivial discrepancy at the critical scale, and is asymptotically distinguishable above it. A square-root accumulation rule would miss this transition because the supports move. The even-contact restriction belongs to this sharp positive-offset benchmark; it is not imposed on the general forward or relative theorems.

A second contribution concerns reconstruction from genuinely unlabelled counts. In a physical three-channel family, the gap and three leading count amplitudes identify the unordered curvature triple and the area, including at coincident curvatures. An ambient formulation with an independent area normalizer uses four amplitudes. A count-only procedure combines a fine timing pilot with fixed-order extrapolation of Bernoulli success probabilities. At extrapolation order $m$, it attains matching curvature error $\varepsilon$ using at most
\begin{equation}\label{eq:v7-intro-count-cost}
 C_m\varepsilon^{-(6+6/m)}\log(C_m/\eta)
\end{equation}
preparations in expectation and with the stated high probability, from a supplied coarse onset bracket.

We prove a matching lower bound for the experiment actually used by this procedure. When each query retains only the indicator of the required collision count, the flight number is bounded, and every actual offset is at most $H_mh$, uniform confidence $1-\eta$ requires at least
\[
       c_m h^{-2}\varepsilon^{-6}\log(1/\eta)
\]
expected preparations. The argument allows adaptive queries and stopping. It uses a rotation-invariant physical family whose opposite perturbations are separated by order $\varepsilon$ in matching curvature distance but by only order $d^2\varepsilon^3$ in finite-offset count probability. Setting $h\asymp\varepsilon^{3/m}$ matches both the accuracy power and the confidence logarithm in \eqref{eq:v7-intro-count-cost}. This is an optimal-order statement within the specified Bernoulli design, not a bound over all collision observations or all schedules.

The inverse results clarify the information being used. In the identical-even class, the one-flight nonlinear germ already determines the contact jets. The finite-flight jet inverse remains uniformly conditioned in flight number at every fixed jet order. Long bridges therefore do not create additional identifiable contact coordinates in that class. Their role is to provide a uniform nonlinear limiting experiment with a quantified sample-size range. Similarly, count-only curvature acquisition uses only $J=3$ or $4$ flight numbers and does not require $j\to\infty$. These two uses of the relative theory should not be combined into an unstated long-record reconstruction claim. In contrast to scalar count amplitudes, two selected-position variances recover both facing curvatures with a Lipschitz inverse even when those curvatures coincide. That improvement comes from a different observation, not from regularizing the same count datum.

\subsection{Observation, information, and asymptotics}
Table~\ref{tab:v7-experiments} records the four experiments in one place. In the table, $\eta$ is a failure probability, $m$ is a fixed extrapolation order, and $\varepsilon$ is a target curvature error. All probability experiments use independent full-phase preparations and noiseless recording. Supplied coarse brackets are not counted as acquired data; fine calibration is charged where stated. Constants may depend on the fixed compact geometric family and on each fixed derivative or extrapolation order.

\begingroup
\setlength{\tabcolsep}{4pt}
\begin{longtable}{@{}p{.16\textwidth}p{.37\textwidth}p{.40\textwidth}@{}}
\caption{Experiments and the variables in their guarantees}\label{tab:v7-experiments}\\
\toprule
Experiment & Observation and supplied information & Target, error, cost, and uniformity\\
\midrule
\endfirsthead
\multicolumn{3}{l}{Table~\thetable\ (continued)}\\
\toprule
Experiment & Observation and supplied information & Target, error, cost, and uniformity\\
\midrule
\endhead
\bottomrule
\endfoot
Long bridge & Endpoints and first residual time, with a failure symbol for a raw preparation. A fixed channel is selected by prescribed patches. Fixed-excess scheduling needs the onset supplied or calibrated. & Geometric parameter unknown. Total variation controls bounded decision risks. Conditional error $Ck\tau^j$; raw error $Cnp_{j,d}\tau^j$. Acquiring $k$ successes costs order $k d^{-2}e^{j\gamma}$. Bounds uniform in $j,d$ on the general compact family. The sharp regime \eqref{eq:v7-intro-phase} adds identical even contacts and $d=o(e^{-j\gamma})$.\\[1ex]
Fixed jets & Coefficients of a normalized nonlinear germ, with its leading gap and curvature specified; identical even facing contacts. This is a coefficient observation, not a sequence of raw counts. & A fixed finite contact jet. Deterministic coefficient error gives Lipschitz jet error with constant $C_M$ at order $M$. Uniform over finite flight numbers and the boundary limit for fixed $M$, not uniformly as $M\to\infty$. No sampling rate is attached to exact germ data.\\[1ex]
Unlabelled Bernoulli counts & Only $\one_{\{N_t\ge j+1\}}$ at each preparation; no channel label or position. A coarse bracket $|g-g_0|\le h/4$ is supplied; $J=3$ in the constrained physical family or $J=4$ with an independent normalizer. & Unknown gap, area, and unordered curvatures. Matching error $C_mh^{m/3}$; cost $C_mh^{-2m-2}\log(C_m/\eta)$ includes fine calibration. Fixed $m,J$, with $h\downarrow0$. The lower bound applies to bounded-$J$, all-history offsets at most $H_mh$ and Bernoulli observations only.\\[1ex]
Selected positions & Successful endpoints on prescribed channel patches, plus the calibration success indicators. Odd $j$; supplied coarse bracket $|j(g-g_0)|\le h/4$. No supplied area or contact frame is used by the two physical variances. & Unknown gap and ordered facing curvatures. Curvature error $C_mh^m$; gap error $C_mh^{m+1}/j$. Cost $C_me^{j\gamma}h^{-2m-2}\log(C_m/\eta)$ includes fine calibration. The inverse constants are uniform in odd $j$ at fixed $m$; preparation cost retains its exponential factor.\\
\end{longtable}
\endgroup

For full embedded collision arrays, our comparison is in the bounded-Lipschitz topology stated in Theorem~\ref{thm:v4-law}, not in total variation. For deterministic physical schedules, the excess is evaluated at the parameter's actual onset; a change of coordinates in a proof is not an observation of an unknown parameter. For every compact minimum-discrepancy reconstruction below, Proposition~\ref{prop:v7-measurable} specifies an exact Borel measurable tie-breaking convention. Existence of a fit therefore defines a statistical decision rule without an additional uniqueness assumption.

\subsection{Relation to existing approaches}
The action-Hessian determinant structure has antecedents in Hill's formula \citep{Hill}. The additional issue here is relative control of a nonlinear twist that vanishes exponentially in the bridge length, followed by the physical residual-time integration. Inverse spectral and marked-length results recover geometric information from different observations and under their respective symmetry and nondegeneracy assumptions \citep{Zelditch,DSKL,FinamoreLeguil}. We do not identify those observations with positive phase-volume onset probabilities. Likewise, confluent Prony analysis provides a relevant framework for coalescing inverse coordinates \citep{BYProny}, but it does not supply the physical timing or rare-event cost of measuring our amplitudes. The detailed comparisons and the complete geometric, circular, marked-response, and fixed-window arguments remain part of the analysis below.

The exact tangent-overlap benchmark and the constant-confidence Bernoulli obstruction were first recorded in the source-pinned referee memorandum \citep{RefereeV6}. We give complete proofs and identify that provenance. The positive-offset joint limit, the high-confidence stopped-design bound, and the measurable reconstruction convention provide the additional links needed for the statistical statements of this article. Neither a diagnostic calculation nor the existence of a reproducible build substitutes for these proofs.
